\chapter{总结与展望}
\label{ch:concl}

\section{工作总结}

本文以 Bender 等人提出的最优时间/空间权衡框架\cite{bender2022otsh} 为理论基础，完成了分层哈希表的结构与算法设计，并据此实现简化系统、开展性能实验。主要难点在于将 MiniArray、k-kick 与 Local Query Router 联调为统一原型，并在 Facility--Cubby--Slot 分层组织下保持辅助结构同步不变量。

在设计层面，第二、三章给出了 Facility--Cubby--Slot 分层组织、MiniArray、Cubby 内 k-kick、Local Query Router 和商压缩等模块的实现方案。MiniArray 以 B 树与 Rank/Select 维护变长元数据；k-kick 在逻辑 bin 间通过有限次踢出完成插入；Local Query Router 以前缀树将查询映射到目标槽位；storage、MiniArray \(B\)、\(M\) 与 Router 在每次搬移后同步更新。

实验测量与配置见第四章表~\ref{tab:exp-n-fixed}--\ref{tab:exp-tier-1e5} 及图~\ref{fig:exp-n-latency}；kick 深度变化时踢出链的最大次数仍不超过对应的 \(k\)，与第二章交换上界一致。

综上，本文完成了从理论结构到可运行原型的设计与实现，记录了工程简化相对严格模型的差异，并给出可复现的测量数据与实验参数。

\section{不足与局限}

受时间与实验条件限制，本文仍存在以下不足：
\begin{itemize}
  \item 未与 IcebergHT\cite{pandey2023iceberght}、PaCHash\cite{kurpicz2023pachash} 等同类系统进行同机对比，尚难从相对角度评价本原型的工程竞争力；
  \item 插入与删除延迟高于查询（\(n=10^5\) 时分别为 60.8\,$\mu$s、2051.9\,$\mu$s，查询 5.1\,$\mu$s），与 MiniArray 维护、k-kick 踢出及 tier 重构等热路径相关；
  \item 第四章合成的 workload 存在分布理想化、阶段化访问模式、规模与场景覆盖窄、缺少外部基线与统计稳健性等局限；
  \item 实验 workload 对键偏斜、大规模数据及 active/old 双表渐进迁移等场景覆盖不足，结论外推范围受测量方式的限制。
\end{itemize}

\section{展望}

本文在 Bender 等人的紧凑哈希框架下完成简化实现与四组测量，归纳工程简化相对理论模型的差异；后续工作可从以下几方面展开：在更大规模与偏斜负载下观察 tier 重建与双表迁移对尾延迟的影响；对插入、删除热路径做针对性优化；选取一至两个紧凑哈希开源实现进行同机对比；补充进程级内存与吞吐的端到端测量，更全面地评估空间--时间权衡\cite{kurpicz2023pachash,lehmann2026mphsurvey}。
